% ===== CLEAN OVERRIDE TABLES =====

\paragraph{PSPEC Proses 3.4: Generate Kartu Ujian}~\\
\begin{longtable}{|p{0.2\textwidth}|p{0.7\textwidth}|}
	\caption{PSPEC Proses 3.4 -- Generate Kartu Ujian}
	\label{table:3.pspec34}\\
	\hline
	\textbf{Nama Proses} & Generate Kartu Ujian \\
	\hline
	\textbf{Deskripsi Proses} & Sistem secara otomatis menggenerasi kartu ujian setelah CAMABA berhasil mengisi dan menyimpan formulir pendaftaran. \\
	\hline
	\textbf{Data Input} & Data CAMABA yang telah terverifikasi dari Data Formulir dan Berkas. \\
	\hline
	\textbf{Data Output} & Kartu ujian dalam format PDF dengan QR Code, tersimpan di Kartu Ujian dan Nomor Ujian. \\
	\hline
	\textbf{Body / Keterangan} & Kartu ujian berisi informasi CAMABA, jadwal ujian, dan QR Code untuk absensi. Kartu dapat diunduh dan dicetak oleh CAMABA. \\
	\hline
\end{longtable}

\paragraph{PSPEC Proses 3.5: Validasi Nomor Ujian}~\\
\begin{longtable}{|p{0.2\textwidth}|p{0.7\textwidth}|}
	\caption{PSPEC Proses 3.5 -- Validasi Nomor Ujian}
	\label{table:3.pspec35}\\
	\hline
	\textbf{Nama Proses} & Validasi Nomor Ujian \\
	\hline
	\textbf{Deskripsi Proses} & Sistem memvalidasi nomor ujian yang diinputkan oleh calon mahasiswa baru (CAMABA) untuk memastikan nomor tersebut terdaftar dan valid. \\
	\hline
	\textbf{Data Input} & Nomor Ujian dari CAMABA. \\
	\hline
	\textbf{Data Output} & Notifikasi "Ditemukan No Pendaftaran" jika nomor valid atau "Tidak Ditemukan No Pendaftaran" jika tidak valid. \\
	\hline
	\textbf{Body / Keterangan} & Proses ini krusial untuk memastikan bahwa hanya CAMABA yang memiliki nomor ujian yang sah yang dapat melanjutkan ke tahap ujian. Jika nomor ujian tidak ditemukan, CAMABA akan menerima pesan kesalahan dan tidak dapat melanjutkan. \\
	\hline
\end{longtable}

\paragraph{PSPEC Proses 4.4: Akses Google Form Ujian}~\\
\begin{longtable}{|p{0.2\textwidth}|p{0.7\textwidth}|}
	\caption{PSPEC Proses 4.4 -- Akses Google Form Ujian}
	\label{table:3.pspec44}\\
	\hline
	\textbf{Nama Proses} & Akses Google Form Ujian \\
	\hline
	\textbf{Deskripsi Proses} & Sistem mengarahkan CAMABA ke Google Form yang berisi soal ujian. Nomor ujian CAMABA akan otomatis terisi (\textbackslash{}textit{prefilled}) di Google Form. \\
	\hline
	\textbf{Data Input} & Status "Waktu Tepat" dari proses 4.2 Validasi Tanggal Ujian. \\
	\hline
	\textbf{Data Output} & Halaman Google Form Ujian dengan nomor ujian yang sudah terisi otomatis. \\
	\hline
	\textbf{Body / Keterangan} & Hal ini memudahkan CAMABA dalam mengikuti ujian karena tidak perlu menginput ulang nomor ujian, serta menjamin integritas data ujian dengan menghubungkan langsung ke nomor ujian yang valid. \\
	\hline
\end{longtable}

\paragraph{PSPEC Proses 6.5: Jawab Pertanyaan (oleh Admin Validasi)}~\\
\begin{longtable}{|p{0.2\textwidth}|p{0.7\textwidth}|}
	\caption{PSPEC Proses 6.5 -- Jawab Pertanyaan (oleh Admin Validasi)}
	\label{table:3.pspec65}\\
	\hline
	\textbf{Nama Proses} & Jawab Pertanyaan (oleh Admin Validasi) \\
	\hline
	\textbf{Deskripsi Proses} & Admin Validasi dapat menjawab pertanyaan yang diajukan oleh CAMABA. \\
	\hline
	\textbf{Data Input} & Jawaban dari Admin Validasi. \\
	\hline
	\textbf{Data Output} & Jawaban terkirim ke CAMABA. \\
	\hline
	\textbf{Body / Keterangan} & Adanya kotak masuk pesan dan tanda pesan yang belum terjawab membantu admin dalam melacak dan menanggapi pertanyaan CAMABA secara efisien. \\
	\hline
\end{longtable}

\paragraph{PSPEC Proses 7.5: Kelola Info Landing Page}~\\
\begin{longtable}{|p{0.2\textwidth}|p{0.7\textwidth}|}
	\caption{PSPEC Proses 7.5 -- Kelola Info Landing Page}
	\label{table:3.pspec75}\\
	\hline
	\textbf{Nama Proses} & Kelola Info Landing Page \\
	\hline
	\textbf{Deskripsi Proses} & Admin Pusat mengelola konten informasi umum yang ditampilkan di halaman utama (\textbackslash{}textit{landing page}) sistem PMB, seperti jadwal, pengumuman, dll. \\
	\hline
	\textbf{Data Input} & Konten informasi yang akan ditampilkan. \\
	\hline
	\textbf{Data Output} & Konfigurasi informasi \textbackslash{}textit{landing page} tersimpan di Konfigurasi Sistem. \\
	\hline
	\textbf{Body / Keterangan} & Memungkinkan admin untuk selalu memperbarui informasi penting bagi CAMABA dan pengunjung. \\
	\hline
\end{longtable}

\paragraph{PSPEC Proses 7.7: Pratinjau Tampilan UI}~\\
\begin{longtable}{|p{0.2\textwidth}|p{0.7\textwidth}|}
	\caption{PSPEC Proses 7.7 -- Pratinjau Tampilan UI}
	\label{table:3.pspec77}\\
	\hline
	\textbf{Nama Proses} & Pratinjau Tampilan UI \\
	\hline
	\textbf{Deskripsi Proses} & Admin Pusat dapat melihat tampilan sistem dari perspektif Admin Validasi atau CAMABA. \\
	\hline
	\textbf{Data Input} & Permintaan pratinjau dari Admin Pusat. \\
	\hline
	\textbf{Data Output} & Tampilan sistem sesuai perspektif yang dipilih. \\
	\hline
	\textbf{Body / Keterangan} & Ini membantu Admin Pusat untuk memverifikasi bagaimana informasi dan perubahan yang mereka buat akan terlihat oleh pengguna lain sebelum dipublikasikan secara luas. \\
	\hline
\end{longtable}
